\section{Related Work}
% \df{we need more related work}
\subsection{Time-Series Forecasting Benchmarks}

A number of large-scale benchmarks evaluate forecasting and time-series models across diverse domains, primarily through numerical prediction or analysis objectives.
TFB~\cite{qiu2024tfb} provides a unified framework for fair comparison of statistical, machine learning, and deep learning forecasters.
GIFT-Eval~\cite{aksu2024gifteval} focuses on standardized evaluation protocols for time-series foundation models under pretraining and finetuning splits, and Chronos~\cite{ansari2024chronos} evaluates forecasting models on a large collection of real-world time series datasets, including both in-domain and zero-shot benchmark settings.
Beyond forecasting, several benchmarks target specific time-series analysis capabilities, such as TSI-Bench~\cite{du2024tsi}, TSB-AD~\cite{liu2024elephant}, and Time-MMD~\cite{liu2024time}, which assess performance on analysis- or distribution-oriented tasks over numerical sequences.
More recent efforts begin to incorporate contextual cues into time-series settings, including CiK~\cite{williams2024context} and TGTSF~\cite{xu2024intervention}, but they still frame evaluation largely as single-task numerical analysis without explicit reasoning-oriented objectives.

Meanwhile, a few benchmarks explore the intersection of language models and numerical time-series by combining question answering with time-series inputs, such as LLM TS Struggle~\cite{merrill2024language}, MTBench~\cite{chen2025mtbench}, and ChatTime~\cite{wang2025chattime}.
While these benchmarks move beyond pure forecasting, they typically do not explicitly incorporate external contextual information or provide a unified multi-task taxonomy to disentangle complementary aspects of temporal competence.
Overall, existing time-series benchmarks are largely optimized for numerical performance, motivating our effort to move beyond prediction error toward understanding- and context-aware temporal reasoning.

\subsection{Temporal Reasoning Benchmarks for LLMs}

A separate line of work evaluates temporal reasoning in large language models using purely textual problems, focusing on symbolic, event-based, and commonsense temporal understanding.
Benchmarks such as TimeBench~\cite{chu2024timebench} and other temporal reasoning evaluations probe abilities such as ordering events, reasoning about duration, and inferring causal-temporal relations from text.
These datasets highlight important deficiencies of LLMs in temporal reasoning, but they omit real numerical time-series signals and do not capture the numerical temporal dynamics central to forecasting and many real-world decision-making settings.

Recent multi-modal benchmarks begin to bridge this gap by coupling numerical time-series with language supervision and downstream reasoning tasks.
For example, TelecomTS~\cite{feng2025telecomts} introduces a large-scale observability benchmark derived from real-world telecommunications networks, supporting tasks such as anomaly detection, root-cause analysis, and question answering over time-series with associated textual context.
Such efforts further motivate the need for benchmarks that jointly evaluate numerical time-series understanding and contextual reasoning, which our work addresses.


\subsection{Agentic and LLM-Based Methods for Time-Series}

Recent work has begun to explore agentic or large language model (LLM)–driven approaches for time-series analysis and forecasting.
TimeSeriesScientist~\cite{zhao2025timeseriesscientist} proposes an agent-based framework that decomposes univariate forecasting into stages such as data inspection, model selection, and result interpretation, demonstrating the potential of agentic workflows for automating parts of the forecasting pipeline.
Other studies investigate the capabilities and limitations of directly applying LLMs to time-series tasks, including numerical forecasting, trend interpretation, and zero-shot reasoning over serialized temporal inputs~\cite{jin2024whatcanllm} .
These works primarily focus on model behavior or system design for specific tasks, rather than providing systematic evaluation across diverse temporal reasoning scenarios.

In parallel, several general-purpose agent frameworks have been proposed to support structured reasoning and coordination among multiple LLM agents.
MetaGPT~\cite{hong2024metagpt} organizes agents according to software engineering roles (e.g., product manager, architect, developer), enabling role-specialized collaboration through explicit communication protocols.
AgentScope~\cite{gao2024agentscope} provides a flexible platform for constructing, orchestrating, and evaluating multi-agent systems, with explicit support for tool use, memory, and interaction logging.
CAMEL~\cite{li2023camel} emphasizes autonomous role-playing and conversational coordination between agents to elicit complex reasoning behaviors without heavy task-specific supervision.
While these frameworks offer powerful abstractions for agent coordination and reasoning, they are largely task-agnostic and do not provide dedicated evaluation settings tailored to time-series understanding.

